% Generated draft: requires mathematical and source review.
\documentclass[11pt]{article}
\usepackage[margin=25mm]{geometry}
\usepackage{fontspec}
\setmainfont{Latin Modern Roman}
\usepackage{amsmath,amssymb,mathtools,unicode-math}
\setmathfont{Latin Modern Math}
\usepackage{enumitem,needspace,xurl,hyperref}
\setlength{\parindent}{0pt}
\setlength{\parskip}{5pt}
\setlength{\emergencystretch}{3em}
\newcommand{\N}{\mathbb{N}}
\newcommand{\Z}{\mathbb{Z}}
\newcommand{\Q}{\mathbb{Q}}
\newcommand{\R}{\mathbb{R}}
\newcommand{\C}{\mathbb{C}}
\newcommand{\F}{\mathbb{F}}
\newcommand{\A}{\mathbb{A}}
\newcommand{\PP}{\mathbb{P}}
\newcommand{\E}{\mathbb{E}}
\newcommand{\eps}{\varepsilon}
\newcommand{\dd}{\,\mathrm{d}}
\begin{document}
\section*{UnsolvedMath ID: 800245}
\textbf{Source problem code:} catalogue entry
\subsection{Definitions and mathematical statement}
The supplied question asks whether finite-type pseudoconvex domains in \(\C^2\) are locally convexifiable by biholomorphic maps having continuous boundary values. The source does not specify boundary regularity, the definition of finite type, or the precise meaning of local convexifiability. The following is a precise formulation conditional on the conventions stated here; these conventions require confirmation against the intended source meaning.

A domain is a connected open subset of \(\C^2\). Call a domain pseudoconvex if it admits a continuous plurisubharmonic exhaustion function. For the conditional formulation, suppose that \(D\subset\C^2\) has \(C^\infty\) boundary. At \(p\in\partial D\), let \(\rho\) be a smooth local defining function, with \(D=\{\rho<0\}\) locally and \(\dd\rho\ne0\) on the boundary. Interpret finite type as finite D'Angelo type:
\[
\Delta_1(\partial D,p)
=\sup_{\gamma}\frac{\operatorname{ord}_0(\rho\circ\gamma)}{\operatorname{ord}_0(\gamma-p)}<\infty,
\]
where the supremum ranges over nonconstant holomorphic curve germs \(\gamma:(\C,0)\to(\C^2,p)\). Here the numerator is the vanishing order of the smooth germ in the real variables underlying \(\C\), and the denominator is the minimum vanishing order of the component germs of \(\gamma-p\); infinite vanishing order is allowed in the numerator.

Under these conventions, the question is whether, for every pseudoconvex domain \(D\subset\C^2\) satisfying
\[
\Delta_1(\partial D,p)<\infty\qquad\text{for every }p\in\partial D,
\]
and every \(p\in\partial D\), there exist an open neighborhood \(U\subset\C^2\) of \(p\), a convex domain \(G\subset\C^2\), and a biholomorphism
\[
F:D\cap U\longrightarrow G
\]
that admits a continuous extension
\[
\overline F:\overline D\cap U\longrightarrow\C^2,
\qquad \overline F|_{D\cap U}=F.
\]
Convexity means that \((1-t)z+tw\in G\) for all \(z,w\in G\) and \(t\in[0,1]\). This formulation interprets local convexifiability as convexity of the entire image of a suitably chosen local piece. The source does not explicitly distinguish this interpretation from convexity only near the image boundary point. Continuous boundary values are interpreted as continuity on the indicated relative closure; injectivity on the boundary and continuous extension of the inverse are not stipulated. No global boundedness assumption is supplied.
\subsection{Short English statement}
The question asks whether every boundary point of a finite-type pseudoconvex domain in \(\C^2\) admits local biholomorphic coordinates making the domain convex, with the coordinate map extending continuously to the boundary. The displayed formulation uses smooth boundary and D'Angelo finite type as explicit, provisional conventions; the source leaves these conventions and the exact local convexification requirement unspecified.
\subsection{Sources}
\begin{itemize}
\item Supplied local UnsolvedMath registry record, permanent ID \(800245\), source problem code \texttt{catalogue entry}, titled ``Several Complex Variables Problem List --- Problem 43.'' The registry attributes the problem to Buzzard, Fornaess, and Krantz, with proposed year 1995, and labels it \texttt{open}. This is a supplied status label, not an independently verified mathematical or literature assessment.
\item Buzzard, Fornaess, and Krantz, \emph{The Several Complex Variables Problem List} (1995), Problem 43. Supplied source URL: \url{https://arxiv.org/abs/math/9509202}. The provenance records access on 2026-10-16 and extraction method \texttt{arxiv\_latex}; these are supplied provenance statements, not a claim of independent access or verification.
\end{itemize}
\end{document}
